\subsection{Ошибка сборки из-за неправильных флагов компилятора}

При попытке собрать Release версию проекта возникла ошибка компиляции.
В файле \texttt{CMakeLists.txt} были указаны флаги оптимизации для компилятора
MSVC (\texttt{/O2}, \texttt{/DNDEBUG}), но проект собирался с помощью компилятора
GCC/MinGW, который использует другой синтаксис флагов.

Рассмотреть ошибку можно ниже.

\begin{lstlisting}[style=CodeListing]
c++.exe: warning: /O2: linker input file unused because linking not done
c++.exe: error: /O2: linker input file not found: No such file or directory
c++.exe: warning: /DNDEBUG: linker input file unused because linking not done
c++.exe: error: /DNDEBUG: linker input file not found: No such file or directory
\end{lstlisting}

Для решения проблемы было необходимо добавить проверку типа компилятора в
\texttt{CMakeLists.txt} и применять соответствующие флаги в зависимости от
используемого компилятора. Для MSVC используются флаги \texttt{/O2} и \texttt{/DNDEBUG}, а
для GCC/MinGW - флаги \texttt{-O2}, \texttt{-DNDEBUG} и \texttt{-s} (для удаления символов).

\begin{lstlisting}[style=CodeListing]
# Release build optimizations
if(CMAKE_BUILD_TYPE STREQUAL "Release")
    if(MSVC)
        # MSVC compiler flags
        set(CMAKE_CXX_FLAGS_RELEASE "${CMAKE_CXX_FLAGS_RELEASE} /O2 /DNDEBUG")
        set(CMAKE_EXE_LINKER_FLAGS_RELEASE "${CMAKE_EXE_LINKER_FLAGS_RELEASE} /OPT:REF /OPT:ICF")
    else()
        # GCC/MinGW compiler flags
        set(CMAKE_CXX_FLAGS_RELEASE "${CMAKE_CXX_FLAGS_RELEASE} -O2 -DNDEBUG")
        set(CMAKE_EXE_LINKER_FLAGS_RELEASE "${CMAKE_EXE_LINKER_FLAGS_RELEASE} -s")
    endif()
endif()
\end{lstlisting}

После внесения изменений в \texttt{CMakeLists.txt} Release версия проекта успешно собирается
как с компилятором MSVC, так и с GCC/MinGW. Флаги оптимизации применяются корректно
в зависимости от используемого компилятора, что обеспечивает оптимальный размер и
производительность исполняемого файла. Проект теперь может быть собран в Release
режиме на различных системах сборки без ошибок компиляции.